\documentclass[11pt,a4paper]{article}
\usepackage[T1]{fontenc}
\usepackage[utf8]{inputenc}
\usepackage{lmodern}
\usepackage{amsmath,amssymb}
\usepackage[margin=25mm]{geometry}
\usepackage[hidelinks]{hyperref}
\hypersetup{pdftitle={What survives refinement: results on Newton's vanishing sagitta and the halving of lattice},pdfauthor={navstokgap research working notes}}
\newcommand{\tightlist}{\setlength{\itemsep}{0pt}\setlength{\parskip}{0pt}}
\setlength{\emergencystretch}{3em}
\setcounter{secnumdepth}{0}
\begin{document}
\hypertarget{what-survives-refinement-results-on-newtons-vanishing-sagitta-and-the-halving-of-lattice-gauge-fields}{%
\section{What survives refinement: results on Newton's vanishing sagitta
and the halving of lattice gauge
fields}\label{what-survives-refinement-results-on-newtons-vanishing-sagitta-and-the-halving-of-lattice-gauge-fields}}

\textbf{Abstract.} We collect the theorems obtained in this repository
on two refinement limits. In Newton's comparison of free motion with
motion under a constant force, the action of a cell of duration \(\tau\)
is spent additively by cuts, and a record of the comparison made by
marks at those cuts reaches exactly the spent action divided by the mark
floor \(\kappa\). The best verdict on the comparison is continuous in
the force if and only if \(\kappa>0\), which identifies a positive floor
with Leibniz's law of continuity read on records. For lattice gauge
theories refined by halving, one directional halving factors into series
moves and a mid-plane integral; for \(SU(2)\) in \(1+2\) dimensions one
halving step is compared with a covariant Gaussian reference, uniformly
in the size of the plane, with responses controlled through third order
and large-field clusters shown to be jointly rare. The one-loop
coefficients transfer to \(SU(N)\) with the factor \(N/2\). Each
statement lists its hypotheses; complete proofs are in the notes cited
after each proof sketch. Planck's constant \(\hbar\) and all dimensional
constants are kept explicit.

\hypertarget{newtons-comparison-and-the-cut-measure}{%
\subsection{1. Newton's comparison and the cut
measure}\label{newtons-comparison-and-the-cut-measure}}

\textbf{Definition 1.1 (Galileo cell).} A body of mass \(m\) moves on
\([0,\tau]\) either freely or under a constant force \(F\), from the
same position and velocity. The two positions differ by
\(P(t)=Ft^2/(2m)\). The \emph{cell action} is
\[K_\tau=\frac{F^2\tau^3}{24\,m}.\]

\textbf{Theorem 1.2 (cut measure).} Cutting the cell at the fraction
\(s\in(0,1)\) of its duration spends the action \(3s(1-s)K_\tau\), that
is \(K_\tau=3s(1-s)K_\tau+K_{s\tau}+K_{(1-s)\tau}\). For any finite set
of cuts with pieces \(\tau_i\),
\[K_\tau=\sum_{\rm cuts}\frac m2\int\dot\phi_{\rm cut}^2+\sum_iK_{\tau_i},\]
where \(\phi_{\rm cut}\) is the Schauder hat inserted by the cut. The
spent action tends to \(K_\tau\) if and only if \(\max_i\tau_i\to0\).

\emph{Proof (sketch).} The hat of a cut at \(s\tau\) has height
\(h=F s(1-s)\tau^2/(2m)\) and energy
\(\frac m2h^2(\frac1{s\tau}+\frac1{(1-s)\tau})=3s(1-s)K_\tau\); later
hats are orthogonal to earlier ones in the Cameron--Martin space
\(H^1_0\), and \(\sum_iK_{\tau_i}\le K_\tau\max_i\tau_i^2/\tau^2\).
\href{cut-measure-newton.md}{Cut-measure note}, Theorems 1--2 and
Proposition 4. \(\square\)

\textbf{Corollary 1.3.} Halving every piece \(n\) times spends
\((1-4^{-n})K_\tau\). Halving only the remaining piece, day after day,
spends \(\frac67(1-8^{-n})K_\tau\to\frac67K_\tau\); two quarter-cuts a
day spend \(\frac{27}{28}K_\tau\).

\hypertarget{records-and-the-floor}{%
\subsection{2. Records and the floor}\label{records-and-the-floor}}

\textbf{Definition 2.1 (mark model).} A protocol places marks at times
\(t_0<t_1<\dots<t_N\) fixed in advance. Mark \(j\) returns the position
with an independent centred Gaussian error of standard deviation
\(\delta_j\) and delivers an independent centred Gaussian impulse of
standard deviation \(\Delta_j\), with
\[\delta_j\Delta_j\ \ge\ \kappa\ \ge0 .\] The two hypotheses are read by
tests invariant under the unknown initial position and velocity. The
optimal such test errs with probability \(\Phi(-d/2)\), where \(d\) is
the statistical distance of the two hypotheses.

\textbf{Theorem 2.2 (the floor).} For every protocol of Definition 2.1,
\(d^2\le K_\tau/\kappa\), independently of the number and times of the
marks; the constant is sharp. Deciding at error \(\epsilon\) therefore
requires \(F^2\tau^3\ge96\,z_{1-\epsilon}^2\,m\kappa\), with
\(z_{1-\epsilon}=\Phi^{-1}(1-\epsilon)\).

\emph{Proof (sketch).} Write the test weights as a piecewise-linear
\(T\) vanishing outside \((t_0,t_N)\); the signal is \(\frac Fm\int T\)
and the noise is at least \(\frac{2\kappa}m\int T'^2\) by the
arithmetic--geometric mean; the sharp inequality
\((\int T)^2\le\frac{\tau^3}{12}\int T'^2\) on \(H^1_0\) closes the
bound. \href{planck-gap-paper.md}{Planck paper}, Theorem 2. \(\square\)

\textbf{Theorem 2.3 (records spend the cut measure).} Assume marks with
\(\delta\Delta=\kappa>0\) are available at every \(\delta>0\). For marks
at \(t_0<\dots<t_N\), with \(\tau=t_N-t_0\) and pieces \(\tau_i\),
\[\sup d^2=\frac{K_\tau-\sum_iK_{\tau_i}}{\kappa},\] approached with
sharp end marks and balanced interior marks. A single cut at the
fraction \(s\) carries \(3s(1-s)K_\tau/\kappa\), and adding a mark
inside a piece adds exactly the share of that cut.

\emph{Proof (sketch).} The maximum of \((\int T)^2/\int T'^2\) over
piecewise-linear \(T\) on the grid is attained by the Galerkin solution
of \(-T''=1\), which is exact at the nodes; its integral is the
trapezoid rule for \(\frac12(w-t_0)(t_N-w)\), whose error is
\(\sum_i\tau_i^3/12\). \href{cut-measure-newton.md}{Cut-measure note},
Proposition 7. \(\square\)

\textbf{Theorem 2.4 (any instrument, bounded aperture).} Let instruments
with any Kraus operators, any apparatus memory and arbitrary adaptivity
act at times in \((0,\tau]\), identical under the two hypotheses,
followed by any final measurement, and let every conditional state keep
the body within position and momentum spreads \(L\) and \(P\). If the
record decides at error \(\epsilon\), then
\[\frac{F\tau L}{\hbar}+\frac{F\tau^2P}{2m\hbar}\ \ge\ 1-2\epsilon .\]
\href{planck-gap-probabilistic.md}{Probabilistic note}, Theorem 2
(hybrid argument over displacement operators).

\textbf{Theorem 2.5 (recoil).} For every protocol of momentum-transfer
marks with pointers read, with probes in arbitrary states, deciding at
error \(\epsilon\) requires
\(s\sum_j\Delta_j\ge8\hbar\arcsin(1-2\epsilon)\), with
\(s=F\tau^2/(2m)\) the sagitta and \(\Delta_j\) the spread of the
impulse delivered by mark \(j\). \href{record-costs-recoil.md}{Recoil
note}, Theorem R, with the constant of the path-length note.

\hypertarget{continuity-of-the-verdict}{%
\subsection{3. Continuity of the
verdict}\label{continuity-of-the-verdict}}

\textbf{Theorem 3.1 (continuity).} In the mark model, the best error
with which the comparison is decided is
\[P_*(F)=\Phi\Bigl(-\tfrac12\sqrt{K_\tau/\kappa}\Bigr)\quad(\kappa>0),\qquad
P_*(F)=0\ \ (F\ne0),\ \ P_*(0)=\tfrac12\quad(\kappa=0).\] Hence \(P_*\)
is continuous at \(F=0\) if and only if \(\kappa>0\). The same holds for
rest against uniform motion with \(mv^2\tau/(8\kappa)\) in place of
\(K_\tau/(4\kappa)\). For a static taper, read without back-action,
\(P_*\) jumps at zero for every \(\kappa\ge0\).

\emph{Proof (sketch).} Theorem 2.2 and its sharpness give the value for
\(\kappa>0\); for \(\kappa=0\) three marks with zero impulse spread and
vanishing resolution decide any \(F\ne0\).
\href{leibniz-continuity-records.md}{Continuity note}, Proposition L.
\(\square\)

\textbf{Corollary 3.2 (topologies).} The supremal statistical distance
between the recorded motions of a cell under forces \(F,F'\) is
\(d(F,F')^2=(F-F')^2\tau^3/(24m\kappa)\) for \(\kappa>0\) and infinite
off the diagonal for \(\kappa=0\). Thus \(\kappa>0\) if and only if
records and the geometry of the motions induce the same topology on the
constant-force family.

\textbf{Corollary 3.3 (known preparations).} Under the hypotheses of
Theorem 2.4, \(P_*(F)\ge\frac12(1-F\tau L/\hbar-F\tau^2P/(2m\hbar))\),
so at bounded aperture the best verdict is continuous at \(F=0\)
whenever \(\hbar>0\).

\textbf{Remark 3.4 (Leibniz, 1687).} Leibniz's law of continuity
requires that when two cases approach, their outcomes approach
(\emph{Nouvelles de la République des Lettres}, July 1687; Gerhardt III,
52). Theorem 3.1 shows that, with outcomes measured by how well they can
be told apart, the law holds for Newton's comparison exactly when
\(\kappa>0\).

\textbf{Proposition 3.5 (the floor as a path-measure threshold).} For
two Gaussian bridges of variance rate \(\hbar/m\) around the chord and
the parabola, the optimal equal-prior test errs with probability
\(\Phi(-\sqrt{K_\tau/(2\hbar)})\), and decides at error \(\epsilon\) iff
\(K_\tau\ge2z_{1-\epsilon}^2\hbar\).
\href{cut-measure-newton.md}{Cut-measure note}, Proposition 6.

\textbf{Proposition 3.6 (one constant).} Brownian free motion with
continuous stationary independent increments and isotropic centred
noise, mass-only laws and independent centre-of-mass composition for all
positive masses force \(mD=\kappa\) with a single constant \(\kappa\)
for all bodies.
\href{stochastic-route-velocitas-ultima.md}{Stochastic-route note},
S1--S3.

\hypertarget{halving-lattice-gauge-fields}{%
\subsection{4. Halving lattice gauge
fields}\label{halving-lattice-gauge-fields}}

Consider a heat-kernel lattice gauge theory in \(D\) dimensions with
compact gauge group \(G\), lattice spacing \(a\) and heat time
\(t=\lambda_Da\) per plaquette.

\textbf{Theorem 4.1 (series/parallel factorization).} One halving of a
lattice direction is exactly the product of series moves, heat-kernel
convolutions that close in every dimension, and a parallel factor
\(\Psi\), a \((D-1)\)-dimensional gauge theory on the mid-plane whose
edge laws are Brownian bridges. It has \(\binom{D-1}2\) transverse
planes. \href{series-parallel-gauge-refinement.md}{Series/parallel
note}, Proposition 1.

\textbf{Theorem 4.2 (free defect).} For the free field the defect
\(\mathcal D\) of one halving is an exact nonpositive quadratic form of
relative size at most \(a^2K^2/16+(Ga)^2/8\); for a cut at the fraction
\(s\) with trapezoid weights the bound scales exactly by \(4s(1-s)\).
\href{series-parallel-gauge-refinement.md}{Series/parallel note},
Proposition 2 and Corollary 2\(_s\).

\textbf{Theorem 4.3 (\(U(1)\) in three dimensions).} One unperturbed
halving step equals the free step up to an extensive error of density
\(e^{-\pi^2/(8\lambda_3a)}\) on the small-field set, for every cut
fraction. The whole measure lies within total variation
\(2(L/a)^3e^{-\pi^2/(6\lambda_3a)}/(1-e^{-\pi^2/(2\lambda_3a)})\) of its
monopole-free part.
\href{series-parallel-gauge-refinement.md}{Series/parallel note},
Theorem 5 and Corollary 5\(_s\);
\href{villain-monopole-refinement.md}{monopole note}, Theorem 2.

\textbf{Theorem 4.4 (bridge midpoints and the centre).} For a compact
connected simply connected group \(G\), the character moments of a
Brownian-bridge point at any cut fraction, and at several cuts, are
exact sums over weights and coroot images. At the reduced fraction
\(s=p/q\) the image terms are central iff \((1-s)H\in P^\vee\); for
\(SU(N)\) they reach the subgroup \(\mathbb Z_{\gcd(q,N)}\) of the
centre, so the centre of \(SU(3)\) is reached only when \(3\mid q\).
\href{sun-midpoint-centre.md}{Centre note}, Theorems 1 and 1\('\).

\hypertarget{the-su2-mid-plane-in-12-dimensions}{%
\subsection{\texorpdfstring{5. The \(SU(2)\) mid-plane in \(1+2\)
dimensions}{5. The SU(2) mid-plane in 1+2 dimensions}}\label{the-su2-mid-plane-in-12-dimensions}}

\textbf{Setting 5.1.} Halve the time direction of the heat-kernel
\(SU(2)\) theory in \(1+2\) dimensions with coarse heat time
\(t=\lambda_3a\le t_0\), on a periodic mid-plane with sides at least
three. In the mid-vertex gauge write \(m_e=m_{*e}e^{\xi_e}\). Fix
\(\delta\in(0,1/10)\) and put \(\varepsilon=t^{1/2-\delta}\),
\(\alpha'=\frac12-3\delta\). The boundary data are small: face
logarithms \(|x_p|,|y_p|\le\varepsilon\) in the two layers and cut data
\(|X_e|\le\varepsilon\). The measure carries a convex radial barrier
vanishing for \(|\xi_e|\le2\varepsilon\) and diverging at the chart
radius; the smallness conditions on \(t\) and the chart are those of the
\href{su2-midplane-small-field.md}{small-field note}, (65), (83), (98)
and §22.3. Let \(C_{m_*}\) be the incidence operator with adjoint
transports of the midpoint connection, \(P=(8+C_{m_*}C_{m_*}^*)^{-1}\)
and \(H=4+\frac12C_{m_*}^*C_{m_*}\).

\textbf{Definition 5.2 (covariant reference).}
\[\mathcal D^{\rm cov}=\frac1{2t}\langle x+y,P(x+y)\rangle-\frac{\|x\|^2+\|y\|^2}{8t}
+\frac12\log\frac{\det H}{\det H_1}.\]

\textbf{Theorem 5.3 (one step).} Uniformly in the size of the plane,
\[\Bigl|\mathcal D_s-\mathcal D^{\rm cov}\Bigr|\le C\,t^{\alpha'}\Bigl[\sum_p\Bigl(1+\frac{|x_p|^2+|y_p|^2}t\Bigr)
+\sum_e\Bigl(1+\frac{|X_e|^2}t\Bigr)\Bigr],\] where \(\mathcal D_s\) is
the normalized log-ratio of the barriered mid-plane integral, including
Haar and heat-kernel amplitudes, and \(C\) depends only on local
constants of the compact chart.

\emph{Proof (sketch).} Classical comparison at the saddle, a local bound
on the saddle determinant, uniform saddle-centred moments by integration
by parts with diagonal dominance, the Laplace remainder by
interpolation, and a Gaussian rectangle inequality for the barrier.
\href{su2-midplane-small-field.md}{Small-field note}, (58) and (62).
\(\square\)

\textbf{Theorem 5.4 (responses through third order).} Let
\(r=\mathcal D_s-\mathcal D^{\rm cov}\), and let \(\partial_a\) change
the boundary link \(a\) at velocity \(\varepsilon\) along an admissible
coordinate box. With \(\gamma=\frac12\log(10/7)\) and the tree length
\(\ell\) of the labels, for \(n=1,2,3\),
\[\sup_{a_1}\sum_{a_2,\dots,a_n}e^{\gamma\ell(a_1,\dots,a_n)}\bigl|\partial_{a_1}\cdots\partial_{a_n}r\bigr|\le C_n\,t^{\alpha'},\]
including adjacent and coincident links, uniformly in the plane and in
the barrier approximation.

\emph{Proof (sketch).} Helffer--Sjöstrand representations of the
covariances, a Feynman--Kac representation of the resolvent with
pointwise block bounds, the differentiated resolvent equation, the
barrier removed one edge at a time, gradient extensions of the scores,
and the cubic vanishing of the local energy difference \(S-V_2\).
\href{su2-midplane-small-field.md}{Small-field note}, (89), (93), (119).
\(\square\)

\textbf{Theorem 5.5 (large-field clusters).} With \(p=2(7/16)^{k_t}\)
and \(k_t=\lfloor3t^{-2\delta}/8\rfloor-1\), for every set \(H\) of
edges and every choice of the barriered edges,
\[\nu\bigl(|\xi_e|>\eta\ \text{for all }e\in H\bigr)\le p^{|H|/16},\]
and for disjoint sets \(H_1,H_2\) at graph distance \(d\) the insertion
products \(F_H=\prod_{e\in H}(e^{-w(\xi_e)}-1)\) satisfy
\[\bigl|E\,F_{H_1\cup H_2}-E\,F_{H_1}E\,F_{H_2}\bigr|\le C\,p^{(|H_1|+|H_2|)/64}(5/12)^d .\]

\emph{Proof (sketch).} Exponential moments of linear functionals under
the uniformly convex measure, stable under linear tilts, and a Chernoff
bound; the correlation by the Helffer--Sjöstrand covariance decay.
\href{su2-midplane-small-field.md}{Small-field note}, (111) and (115).
\(\square\)

\hypertarget{from-su2-to-sun}{%
\subsection{\texorpdfstring{6. From \(SU(2)\) to
\(SU(N)\)}{6. From SU(2) to SU(N)}}\label{from-su2-to-sun}}

\textbf{Theorem 6.1 (one-loop coefficients).} With the metric
\(\langle X,Y\rangle=-2\operatorname{tr}XY\), the relative-order-\(t\)
coupling coefficients of one halving for \(SU(N)\) are those of
\(SU(2)\) multiplied by \(N/2\):
\[\delta_t^{12}=\delta_t^{13}=\frac{Nt}{48}\int_{\mathcal B}\frac qR,\qquad
\delta_t^{23}=\frac N2\,t\Bigl(-\frac1{12}+2\mathcal I\Bigr),\] with the
group-independent integrals of the
\href{su2-midplane-order-t.md}{order-\(t\) note}. For \(SU(3)\) the cut
couplings shift by \(\frac t{16}\int q/R>0\) and the transverse coupling
by \(t(-\frac18+3\mathcal I)<0\).

\emph{Proof (sketch).} Every group factor enters through the Killing
form, \(\operatorname{tr}_{\rm adj}(\operatorname{ad}X)^2=-N|X|^2\);
around a Cartan background the plaquette Hessian splits into root
spaces, each an isometrically embedded \(SU(2)\) of charge
\(\alpha(Y)\), and \(\sum_{\alpha>0}\alpha(Y)^2=\frac N2|Y|^2\).
\href{su2-midplane-order-t.md}{Order-\(t\) note}, §7b. \(\square\)

\textbf{Proposition 6.2 (twisted \(SU(3)\) sector).} For the clock and
shift matrices \(P,Q\in SU(3)\), \(PQ=\omega QP\) with
\(\omega=e^{2\pi i/3}\), the adjoint action has the eight joint phases
\((2\pi a/3,-2\pi b/3)\), \((a,b)\in\mathbb Z_3^2\setminus\{0\}\). The
twisted flat sector therefore has no adjoint zero mode, and its
mid-plane Laplacian is bounded below by \(2-2\cos(2\pi/(3n_{\max}))\).
\href{su2-midplane-small-field.md}{Small-field note}, §16.

\textbf{Proposition 6.3 (bridge tails for any compact group).} A
bi-invariant metric has nonnegative Ricci curvature, so the Li--Yau
bounds give, for the bridge midpoint of duration \(t\le1\) and every
\(\epsilon'\in(0,1)\),
\(\beta\{d(m,m_*)\ge r\}\le C_{\epsilon'}t^{-n/2}\exp\{-[4(r-|X|/2)^2/(2+\epsilon')-|X|^2/(2-\epsilon')]/t\}\),
with \(n=\dim G\). \href{su2-midplane-small-field.md}{Small-field note},
§16.

\hypertarget{consequence-for-state}{%
\subsection{7. Consequence for STATE}\label{consequence-for-state}}

This note is the formal compendium of the results above, written for
readers who want the statements with their hypotheses in one place; the
cited notes hold the proofs and remain the record. It is linked from the
site's index and tutorial.
\end{document}
